% Chapter 7: Comparison: Describing Differences

\chapter{Comparison: Describing Differences}

\section{Pedagogical Purpose}
This chapter builds on the understanding of adjectives by introducing how they can be used to compare nouns. This is fundamental for nuanced description and expression.

\begin{learningobjectives}
The learner can:
\begin{itemize}
    \item define the positive, comparative, and superlative degrees of adjectives.
    \item form the comparative and superlative degrees of regular adjectives.
    \item identify and use common irregular comparative and superlative adjectives (e.g., good, better, best).
    \item use comparative adjectives to compare two nouns.
    \item use superlative adjectives to compare three or more nouns.
\end{itemize}
\end{learningobjectives}

\section{Prerequisite Check}
\begin{quickcheck}
Underline the adjective in each sentence.
\begin{enumerate}
    \item The \underline{red} car is fast.
    \item She has \underline{long} hair.
    \item There are \underline{five} apples in the basket.
\end{enumerate}
\end{quickcheck}

\section{Chapter Opener}
\subsection*{A Classroom Conversation}
\textbf{Ms. Sharma:} "Look at these three animals: a rabbit, a dog, and an elephant. How would you describe their size?"

\textbf{Rohan:} "The rabbit is \textbf{small}."

\textbf{Maya:} "The dog is \textbf{smaller} than the elephant."

\textbf{Rohan:} "And the elephant is the \textbf{biggest} of all!"

\textbf{Ms. Sharma:} "Excellent! You just used different forms of adjectives to compare the animals. You used 'small' to describe one animal, 'smaller' to compare two animals, and 'biggest' to compare three or more. These are called the \textbf{degrees of comparison} for adjectives."

\section{Concept Discovery (Degrees of Comparison)}
Adjectives can change their form to show comparison. There are three degrees of comparison:

\subsection*{A. Positive Degree}
The \textbf{Positive Degree} of an adjective describes a single noun without making any comparison.
\begin{examplebox}
    The rabbit is \textbf{small}.
\end{examplebox}

\subsection*{B. Comparative Degree}
The \textbf{Comparative Degree} of an adjective is used to compare two nouns. We usually add \textbf{-er} to the adjective or use the word \textbf{more} before it.
\begin{examplebox}
    The dog is \textbf{smaller} than the elephant.
\end{examplebox}

\subsection*{C. Superlative Degree}
The \textbf{Superlative Degree} of an adjective is used to compare three or more nouns. We usually add \textbf{-est} to the adjective or use the word \textbf{most} before it.
\begin{examplebox}
    The elephant is the \textbf{biggest} animal.
\end{examplebox}

\section{Forming Comparative and Superlative Degrees}

\subsection*{A. Regular Adjectives (Adding -er and -est)}
For most short adjectives (one or two syllables), we add \textbf{-er} for the comparative degree and \textbf{-est} for the superlative degree.

\begin{table}[h!]
\centering
\begin{tabular}{|l|l|l|}
\hline
\textbf{Positive} & \textbf{Comparative} & \textbf{Superlative} \\
\hline
tall & taller & tallest \\
fast & faster & fastest \\
kind & kinder & kindest \\
small & smaller & smallest \\
long & longer & longest \\
\hline
\end{tabular}
\end{table}

\subsection*{B. Spelling Rules for -er and -est}
Sometimes, we need to change the spelling of the adjective before adding -er or -est.

\begin{table}[h!]
\centering
\begin{tabular}{|l|l|l|l|}
\hline
\textbf{Rule} & \textbf{Example} & \textbf{Comparative} & \textbf{Superlative} \\
\hline
If the adjective ends in \textbf{-e}, just add \textbf{-r} and \textbf{-st}. & large & larger & largest \\
& wise & wiser & wisest \\
\hline
If the adjective ends in a \textbf{consonant + y}, change \textbf{y} to \textbf{i} and add \textbf{-er} and \textbf{-est}. & happy & happier & happiest \\
& busy & busier & busiest \\
\hline
If the adjective ends in a \textbf{single vowel + single consonant}, double the final consonant. & big & bigger & biggest \\
& hot & hotter & hottest \\
\hline
\end{tabular}
\end{table}

\subsection*{C. Long Adjectives (Using 'more' and 'most')}
For longer adjectives (usually with two or more syllables), we do not add -er or -est. Instead, we use \textbf{more} for the comparative and \textbf{most} for the superlative.

\begin{table}[h!]
\centering
\begin{tabular}{|l|l|l|}
\hline
\textbf{Positive} & \textbf{Comparative} & \textbf{Superlative} \\
\hline
beautiful & more beautiful & most beautiful \\
difficult & more difficult & most difficult \\
important & more important & most important \\
\hline
\end{tabular}
\end{table}

\subsection*{D. Irregular Adjectives}
Some adjectives do not follow any rules. They change completely. You have to remember them!

\begin{table}[h!]
\centering
\begin{tabular}{|l|l|l|}
\hline
\textbf{Positive} & \textbf{Comparative} & \textbf{Superlative} \\
\hline
good & better & best \\
bad & worse & worst \\
little & less & least \\
much / many & more & most \\
far & farther / further & farthest / furthest \\
\hline
\end{tabular}
\end{table}

\begin{rememberbox}
    The \textbf{Comparative} form is often followed by the word \textbf{than}.
    \begin{itemize}
        \item \textit{Example:} Rohan is taller \textbf{than} Maya.
    \end{itemize}
    The \textbf{Superlative} form is often used with the word \textbf{the} before it.
    \begin{itemize}
        \item \textit{Example:} Rohan is \textbf{the} tallest boy in the class.
    \end{itemize}
\end{rememberbox}

\section{Summary}
\begin{summarybox}
    \begin{itemize}
        \item \textbf{Positive Degree:} Describes one thing. (e.g., a \textbf{fast} car)
        \item \textbf{Comparative Degree:} Compares two things. (e.g., a \textbf{faster} car)
        \item \textbf{Superlative Degree:} Compares three or more things. (e.g., the \textbf{fastest} car)
        \item We change adjectives to compare things, using \textbf{-er/-est} or \textbf{more/most}. Some adjectives have special forms!
    \end{itemize}
\end{summarybox}